\documentclass[11pt, a4paper]{article}

% ── PACKAGES ──────────────────────────────────────────────────
\usepackage[utf8]{inputenc}
\usepackage[T1]{fontenc}
\usepackage{geometry}
\geometry{
  top=2.5cm,
  bottom=2.5cm,
  left=2.5cm,
  right=2.5cm
}
\usepackage{hyperref}
\hypersetup{
  colorlinks=true,
  urlcolor=blue,
  linkcolor=blue,
  citecolor=blue,
  pdftitle={Mars First: A Geometric Derivation of
    Martian Biological Priority and the Inversion
    of the Panspermia Prior},
  pdfauthor={Eric Robert Lawson}
}
\usepackage{microtype}
\usepackage{parskip}
\usepackage{booktabs}
\usepackage{array}
\usepackage{tabularx}
\usepackage{xcolor}
\usepackage{mdframed}
\usepackage{amsmath}
\usepackage{amssymb}
\usepackage{enumitem}
\usepackage{titlesec}
\usepackage{lmodern}
\usepackage{authblk}
\usepackage{abstract}

% ── COLOURS ───────────────────────────────────────────────────
\definecolor{headerblue}{RGB}{20,60,140}
\definecolor{boxblue}{RGB}{20,60,140}
\definecolor{rulecolor}{RGB}{20,60,140}

% ── SECTION FORMAT ────────────────────────────────────────────
\titleformat{\section}
  {\large\bfseries\color{headerblue}}
  {\thesection.}{0.5em}{}
\titleformat{\subsection}
  {\normalsize\bfseries\color{headerblue}}
  {\thesubsection.}{0.5em}{}
\titleformat{\subsubsection}
  {\normalsize\itshape\color{headerblue}}
  {\thesubsubsection.}{0.5em}{}

\setlength{\parskip}{0.6em}
\setlength{\parindent}{0pt}

% ══════════════════════════════════════════════════════════════
\begin{document}

% ── TITLE BLOCK ───────────────────────────────────────────────
\title{
  \large\bfseries\color{headerblue}
  MARS FIRST\\[0.4em]
  \normalsize\color{headerblue}
  A Geometric Derivation of Martian Biological
  Priority,\\
  the Inversion of the Panspermia Prior
  Probability,\\
  and the Testable Prediction That the Oldest
  Known Biosphere\\
  in the Solar System Is Fossilised on Mars
}

\author[1]{Eric Robert Lawson}
\affil[1]{OrganismCore (Independent Research)\\
  ORCID: \href{https://orcid.org/0009-0002-0414-6544}
  {0009-0002-0414-6544}\\
  \href{mailto:OrganismCore@proton.me}
  {OrganismCore@proton.me}
}

\date{2026-03-13\\[0.3em]
  \footnotesize
  LECA Pre-registration DOI:
  \href{https://doi.org/10.5281/zenodo.18986790}
  {10.5281/zenodo.18986790}
  (timestamp: 2026-03-12)\\
  Mars Collapse vs.\ Emergence DOI:
  \href{https://doi.org/10.5281/zenodo.18995862}
  {10.5281/zenodo.18995862}
  (timestamp: 2026-03-13)\\
  Universal Life Detection Standard DOI:
  \href{https://doi.org/10.5281/zenodo.18994082}
  {10.5281/zenodo.18994082}\\
  Repository:
  \href{https://github.com/Eric-Robert-Lawson/attractor-oncology}
  {github.com/Eric-Robert-Lawson/attractor-oncology}
}

\maketitle
\thispagestyle{empty}

% ── STATEMENT BOX ───────────────────────────────���─────────────
\begin{mdframed}[
  backgroundcolor=boxblue,
  linecolor=boxblue,
  linewidth=0pt,
  innertopmargin=0.5em,
  innerbottommargin=0.5em,
  innerleftmargin=1em,
  innerrightmargin=1em
]
\begin{center}
\small\bfseries\color{white}
Mars was habitable before Earth was stably
habitable. Mars lost its magnetic field at
$\sim$4.0--4.1\,Ga. The Cheyava Falls biosignature
features (Hurowitz et al., \textit{Nature} 645,
2025) date to $\sim$3.8--4.1\,Ga and are consistent
with a mature, metabolically diverse LUCA-grade
biosphere in terminal collapse --- not emergence.
The earliest confirmed Earth biosignatures date to
$\sim$3.5--3.8\,Ga. The Martian biosphere collapse
therefore occurred at, or prior to, the emergence
of life on Earth. This geometric derivation,
produced 2026-03-13 from first principles in
cancer biology without prior knowledge of
astrobiology, constitutes the pre-registration
of three testable predictions: (1) the Martian
biosphere predates the earliest Earth
biosignatures, (2) the panspermia prior
probability is no longer negligible and is the
default hypothesis requiring elimination rather
than the exotic hypothesis requiring proof, and
(3) the Mars sample return complexity gradient
will show biological activity predating
$\sim$3.8\,Ga. Pre-registered: 2026-03-13.
\end{center}
\end{mdframed}

\vspace{1em}

% ── ABSTRACT ──────────────────────────────────────────────────
\begin{abstract}
\noindent
We present a geometric derivation of Martian
biological priority --- the conclusion that the
Martian biosphere existed before any confirmed
Earth biosignature, and that the collapse of
that biosphere is recorded in the Cheyava Falls
features reported by Hurowitz et al.\
(\textit{Nature} 645, 2025).
The argument proceeds from four independent
lines of evidence: (1) the physics of planetary
cooling, which establishes that Mars reached a
chemically stable, water-bearing, magnetically
shielded surface condition before Earth;
(2) the confirmed geochemical maturity of the
Cheyava Falls mineral assemblage, which indicates
a long-established rather than nascent biological
community; (3) the atmospheric collapse timeline,
which places the Cheyava Falls fossilisation at
the opening of planetary uninhabitability rather
than the establishment of habitability; and
(4) the time-constraint argument, which
establishes that LECA-grade signatures in
pre-GOE conditions are geometrically impossible
under emergence but expected under collapse.
We derive the inversion of the panspermia prior:
given Martian biological priority and confirmed
Mars-to-Earth meteorite transfer, the hypothesis
that Earth life originated from Mars is the
default hypothesis requiring elimination, not
the exotic hypothesis requiring proof.
We pre-register three testable predictions
resolvable by Mars sample return and genomic
analysis of returned material.
This derivation originated in attractor geometry
applied to cancer drug targets and contains
no prior knowledge of astrobiology.
Timestamp: 2026-03-13.
\end{abstract}

\vspace{0.5em}
\noindent\textcolor{rulecolor}{\rule{\linewidth}{0.6pt}}

% ── SECTION 1 ─────────────────────────────────────────────────
\section{Derivation Origin and Epistemic Status}

This document is the fifth in a pre-registration
chain that began on 2026-03-12 with the LECA
Resurrection Protocol
\cite{preregistration}.
The chain proceeded through the LUCA companion
paper \cite{lucapaper}, the Universal Life
Detection Standard \cite{uldstandard}, and the
Mars Collapse vs.\ Emergence pre-registration
\cite{marspreregistration}.

The derivation in each document originated from
attractor geometry applied to cancer drug targets,
not from astrobiology or Mars mission data.
The conclusions are therefore not post-hoc
rationalisations of existing data.
They are geometric predictions that arrived at
astrobiology as a necessary destination.

This document pre-registers the conclusion that
was derived in the session of 2026-03-13:

\textbf{Mars was biologically active before Earth.}
\textbf{The panspermia prior is inverted.}
\textbf{These predictions are testable.}

The honest epistemic status is stated precisely
in Section~7.

% ── SECTION 2 ─────────────────────────────────────────────────
\section{The Planetary Cooling Argument}

\subsection{The Physics}

A planetary body loses internal heat at a rate
inversely proportional to its size.
Smaller bodies have a higher surface-area-to-volume
ratio and radiate heat more efficiently per unit
of internal mass.

Mars has approximately 15\% of Earth's mass and
approximately 28\% of Earth's surface area.
Mars therefore cooled faster than Earth.
This is not a hypothesis.
It is a consequence of geometry and thermodynamics.

\subsection{The Consequence for Habitability}

A cooler planet solidifies its crust earlier,
retains a magnetic dynamo for a shorter period
but establishes stable surface conditions sooner.
The sequence on Mars:

\begin{enumerate}[leftmargin=*, label=(\arabic*)]
\item Crust solidified and stable liquid water
  possible: $\sim$4.5\,Ga or earlier.
\item Magnetic field active and shielding
  surface from solar wind: confirmed present
  at 4.1\,Ga and earlier by crustal
  magnetisation data.
\item Hydrothermal systems from residual
  volcanic activity: confirmed by geology.
\item Iron-rich anoxic ocean conditions:
  confirmed by Cheyava Falls mineral data
  \cite{hurowitz2025}.
\item Phosphate availability 5--10$\times$ higher
  than early Earth: confirmed by Adcock et al.\
  \cite{adcock2013}.
\end{enumerate}

\subsection{Earth's Situation at the Same Time}

Earth at 4.4--4.5\,Ga:

\begin{enumerate}[leftmargin=*, label=(\arabic*)]
\item Late Heavy Bombardment ongoing.
  Surface repeatedly sterilised by impacts.
\item Crust unstable.
  Liquid water present but episodic.
\item No confirmed biosignatures at
  this age on Earth.
\item Hadean eon.
  Named, with geological justification,
  after hell.
\end{enumerate}

The oldest confirmed evidence for liquid
water on Earth comes from Jack Hills zircons
dated to $\sim$4.4\,Ga \cite{wilde2001}.
The oldest confirmed Earth biosignatures date
to $\sim$3.5--3.7\,Ga (Dresser Formation
stromatolites, Pilbara Craton, Australia).
Some carbon isotope evidence suggests possible
biological activity at $\sim$3.8--4.1\,Ga,
but this remains contested.

\textbf{Mars had stable, magnetically shielded,
mineral-rich, phosphate-enriched liquid water
before Earth had any of these conditions
simultaneously.}

The origin of life requires all of these
conditions simultaneously.
Mars satisfied them first.

% ── SECTION 3 ─────────────────────────────────────────────────
\section{The Mature Community Argument}

\subsection{What Cheyava Falls Shows}

The Cheyava Falls mineral assemblage
\cite{hurowitz2025}:

\begin{itemize}[leftmargin=*]
\item Vivianite (Fe$_3$(PO$_4$)$_2$):
  requires coupled iron-reducing and
  phosphate-cycling metabolism.
  Abundant and spatially organised.
\item Greigite (Fe$_3$S$_4$):
  requires sulfate or sulfur-reducing
  metabolism producing H$_2$S.
\item Both minerals co-occurring:
  requires simultaneously active
  iron, sulfur, and phosphate cycling
  in distinct spatial niches.
\item Organic carbon co-located:
  detected by SHERLOC fluorescence
  spectroscopy.
\item Colony-scale spatial organisation:
  leopard spots at 1--2\,mm, poppy seeds
  at $<$200\,$\mu$m.
\end{itemize}

\subsection{What Co-occurrence Requires}

Vivianite and greigite co-occurrence requires
metabolic niche partitioning.
Vivianite forms where Fe$^{2+}$ meets phosphate
with low local sulfide.
Greigite forms where Fe$^{2+}$ meets H$_2$S
from sulfate reduction.
For both to occur simultaneously in adjacent
spatial zones, two distinct microbial communities
must be operating, partitioned by the
microscale redox chemistry they themselves create.

This is not the signature of pioneering
colonisation.
On Earth, metabolic niche partitioning in
iron-sulfur-phosphate systems develops in
mature microbial mat communities that have
been established for extended periods.
Early colonisers produce a single dominant
mineral phase.
Mature communities produce coupled, spatially
organised multi-phase mineral records.

\subsection{The Implication for Community Age}

A mature LUCA-grade community producing the
Cheyava Falls mineral assemblage did not
appear weeks or years before fossilisation.
It required time to establish:

\begin{itemize}[leftmargin=*]
\item Time for iron-reducing communities
  to establish and modify local geochemistry.
\item Time for sulfur-reducing communities
  to occupy the niche created by iron-reducers.
\item Time for phosphate cycling to intensify
  as community density increased.
\item Time for the coupled vivianite-greigite
  mineral record to accumulate at
  colony scale.
\end{itemize}

On Earth, analogous mature microbial mat
communities in iron-sulfur-phosphate systems
require millions to tens of millions of years
to develop from initial colonisation to
the degree of metabolic organisation
visible at Cheyava Falls.

If the Cheyava Falls community was already
mature at 3.8--4.1\,Ga, then it began
significantly earlier.

The plausible origin window for the
Martian biosphere is therefore 4.3--4.5\,Ga ---
during or before the period when Earth was
being repeatedly sterilised by Late Heavy
Bombardment.

% ── SECTION 4 ─────────────────────────────────────────────────
\section{The Time Constraint Argument}

\subsection{The Cascade Duration on Earth}

On Earth:

\begin{itemize}[leftmargin=*]
\item First confirmed life: $\sim$3.5--3.8\,Ga.
\item Great Oxidation Event (GOE): $\sim$2.4\,Ga.
\item Time from first life to GOE:
  $\sim$1.0--1.4 billion years.
\end{itemize}

That interval represents the time required for:
LUCA-grade life to establish;
photosynthesis to evolve;
photosynthetic communities to reach sufficient
density to change atmospheric oxygen;
oxygen to accumulate against geological sinks;
and the GOE to occur.

\subsection{Mars Had Insufficient Time}

Mars had stable surface water conditions
from approximately 4.5\,Ga to approximately
3.0\,Ga, with the most stable conditions in
the Noachian (4.1--3.7\,Ga).
The magnetic field was lost at 4.0--4.1\,Ga.
From that moment, atmospheric stripping
began at rates hundreds of times higher
than today (MAVEN mission data).
By 3.7\,Ga, Mars had lost approximately
90--95\% of its atmosphere.

The window from magnetic field loss to
effective uninhabitability was approximately
300--400 million years.

Earth required 1.0--1.4 billion years
to complete the LUCA-to-GOE cascade.

Mars had less than one-third of that time.

\textbf{The biological cascade from LUCA-grade
life to a GOE-equivalent was geometrically
impossible on Mars given the available time
window.}

This is not a probability statement.
It is a time constraint derived from
the confirmed timeline and the confirmed
duration of the cascade on Earth.

\subsection{The Consequence for
the LECA Signal}

The LECA-grade morphological signature at
Cheyava Falls exists in pre-GOE, State~1
water conditions.
Under emergence, a LECA-grade organism
in State~1 water without a GOE is
a geometric contradiction:
the LECA attractor requires the GOE to
raise Axis~2 (oxygen) above the
destabilisation threshold
\cite{marspreregistration}.

Under collapse, the LECA-grade signature
requires no contradiction.
Complex organisms evolved earlier,
when conditions were better.
They were dying in State~1 water at
the time of fossilisation because
conditions had deteriorated.
The LECA attractor was not stable ---
it was in terminal decline.

The LECA signal at Cheyava Falls
is therefore only coherent under collapse.

% ── SECTION 5 ─────────────────────────────────────────────────
\section{The Inversion of the Panspermia Prior}

\subsection{The Standard Prior}

The standard scientific prior for panspermia
--- the transfer of life between planetary bodies ---
is that it is an exotic hypothesis.

The default assumption is independent origin:
Earth life began on Earth.

Panspermia is the hypothesis that requires
extraordinary evidence to support.

\subsection{Why the Prior Is Inverted
by This Derivation}

The standard prior was set in a framework
where Earth was assumed to be the first
habitable planet in the solar system.

This derivation establishes that Mars
was habitable first.

Given:

\begin{enumerate}[leftmargin=*, label=(\arabic*)]
\item Mars was stably habitable
  before Earth was stably habitable.
\item Mars had life at $\sim$4.4--4.5\,Ga
  (geometric derivation from mature
  community argument and planetary
  cooling physics).
\item Earth life begins at $\sim$3.5--3.8\,Ga
  (confirmed) or possibly earlier
  (contested).
\item Mars-to-Earth meteorite transfer
  is confirmed real:
  we have Martian meteorites on Earth
  (ALH84001 and others).
  Transfer occurs continuously via
  impact ejecta.
\item The transfer window --- from Mars
  life origin at 4.4--4.5\,Ga to
  Earth life first appearing ---
  spans hundreds of millions of years
  during which Mars rocks were
  landing on Earth.
\end{enumerate}

Under these conditions:

\textbf{The hypothesis that Earth life
originated from Mars is not exotic.
It is the default hypothesis that requires
elimination.}

Not because panspermia is proven.
But because the prior conditions that
made independent Earth origin the
default assumption no longer hold.

Mars was first.
Mars had life.
Mars was transferring material to Earth.
Earth life appeared later.

The burden of proof has shifted.

\subsection{What Would Resolve the Question}

The common origin vs.\ independent emergence
distinction is resolvable by genomic analysis
of Mars sample return material.

\begin{itemize}[leftmargin=*]
\item If Martian organisms share
  homologous genetic sequences,
  ribosomal RNA structure, or
  biochemical architecture with
  Earth life: common origin confirmed.
  Either Earth-to-Mars or
  Mars-to-Earth transfer.
  Given the timeline, Mars-to-Earth
  is the more probable direction.
\item If Martian organisms show
  fundamentally different genetic
  architecture with no homology
  to Earth life: independent emergence.
  Life arose independently on two
  planets in the same solar system.
  The LECA attractor is confirmed
  as a universal geometric consequence
  of carbon chemistry.
  Life is not rare in the universe.
\end{itemize}

Both outcomes are scientifically
extraordinary.
Both are testable.
Neither requires Mars sample return
to be pre-registered.
This document pre-registers the
prediction that one of these two
outcomes will obtain.

% ── SECTION 6 ─────────────────────────────────────────────────
\section{The Three Pre-Registered Predictions}

\begin{mdframed}[
  linecolor=rulecolor,
  linewidth=0.8pt,
  innertopmargin=0.6em,
  innerbottommargin=0.6em,
  innerleftmargin=1em,
  innerrightmargin=1em
]
\textbf{Pre-Registered Predictions
(locked 2026-03-13, prior to Mars
sample return and prior to any
genomic analysis of Mars material):}

\medskip
\textbf{Prediction 1 --- Martian Biological Priority:}

The Mars sample return subsurface
sedimentary record from Jezero Crater
will reveal biological activity
(biosignature complexity gradient,
mineral assemblage, organic carbon)
dating to $\geq$3.8\,Ga ---
at or before the earliest confirmed
Earth biosignatures.

\textit{Test:} Radiometric dating of
biosignature-bearing layers in
sample return material.
Comparison against the $\sim$3.5--3.7\,Ga
Dresser Formation stromatolites and
$\sim$3.8\,Ga Isua carbon isotope record.

\medskip
\textbf{Prediction 2 --- Panspermia Default:}

Genomic or biochemical analysis of Mars
sample return material will reveal either:

(a) Homology with Earth biochemistry,
confirming common origin (direction of
transfer to be determined by comparative
phylogenetics and timing), or

(b) Independent biochemical architecture,
confirming independent emergence and
the universality of the LECA attractor.

\textit{There is no third option.}
The prediction is that the question is
resolvable and that the answer will
be one of these two.

\medskip
\textbf{Prediction 3 --- The Vivianite Horizon
as Death Marker:}

The stratigraphic layer immediately above
the Cheyava Falls vivianite-bearing unit
will show a sharp transition to abiotic
iron chemistry (hematite, magnetite,
siderite) with no vivianite and no
co-located organic carbon.

The vivianite horizon marks the
biological cessation boundary of the
Martian surface biosphere.

\textit{Test:} Mineralogical and organic
carbon analysis of material immediately
above and below the Cheyava Falls unit
in the sample return core.
\end{mdframed}

% ── SECTION 7 ─────────────────────────────────────────────────
\section{Honest Statement of Epistemic Status}

\subsection{What Is Geometrically Derived and
Strongly Supported}

\begin{enumerate}[leftmargin=*, label=(\arabic*)]
\item Mars was physically habitable before
  Earth was stably habitable.
  This follows from planetary cooling physics
  and is not disputed.

\item The Cheyava Falls mineral assemblage
  is consistent with a mature, established
  biological community rather than an
  emergent one.
  This follows from the vivianite-greigite
  co-occurrence chemistry and the
  metabolic niche partitioning argument.

\item LECA-grade signatures in pre-GOE
  State~1 water are geometrically
  impossible under emergence and
  expected under collapse.
  This follows from the two-axis
  LECA stabilisation model
  \cite{marspreregistration}.

\item Given Martian biological priority
  and confirmed Mars-to-Earth transfer,
  the panspermia hypothesis is the
  default hypothesis requiring
  elimination, not the exotic one
  requiring proof.
  This follows from the reversal of
  the prior conditions that established
  the original standard prior.
\end{enumerate}

\subsection{What Is Not Yet Confirmed}

\begin{enumerate}[leftmargin=*, label=(\arabic*)]
\item The Cheyava Falls features are
  biological and not abiotic.
  NASA's Step~1 classification reflects
  this uncertainty.
  ARM~A physical anchoring of the LECA
  morphological standard is pending.

\item The Martian biosphere predated
  Earth life specifically by any
  particular margin.
  The timing argument is geometric
  and probabilistic, not precisely dated.

\item The direction of any panspermia
  transfer.
  This requires genomic analysis of
  Mars sample return material.

\item Whether Earth life IS Martian life.
  This is the most extreme form of the
  panspermia conclusion.
  It is not claimed here as confirmed.
  It is claimed as the default hypothesis
  requiring elimination given the
  geometric derivation.
\end{enumerate}

\subsection{What ARM A and Sample Return
Confirm}

\begin{itemize}[leftmargin=*]
\item ARM~A ($\sim$\$200, BSL-1, Rogers Park):
  Physically anchors the LECA morphological
  reference standard.
  Moves Premise~1 from geometrically derived
  to physically confirmed.

\item Mars sample return (2030s):
  Complexity gradient in subsurface sediment.
  Confirms collapse or emergence.
  Dating of deepest biosignature-bearing layers.
  Confirms or denies Martian biological priority.

\item Genomic analysis of sample return material:
  Resolves common origin vs.\ independent emergence.
  Determines panspermia vs.\ convergence.
  Determines direction of transfer if common origin.
\end{itemize}

% ── SECTION 8 ─────────────────────────────────────────────────
\section{The Derivation Chain That Produced
This Conclusion}

\begin{quote}
Machine consciousness question\\
$\rightarrow$ Triadic invariant\\
$\rightarrow$ Identity Axis Theorem\\
$\rightarrow$ Cancer drug target geometry\\
$\rightarrow$ Reproduction Theorem\\
$\rightarrow$ Inverse Accessibility Theorem\\
$\rightarrow$ LECA Resurrection Protocol
  \cite{preregistration}\\
$\rightarrow$ LECA attractor as universal
  base state \cite{uldstandard}\\
$\rightarrow$ LECA Attractor Habitability Standard\\
$\rightarrow$ Application to ancient Mars\\
$\rightarrow$ Cheyava Falls convergence\\
$\rightarrow$ Collapse vs.\ emergence
  pre-registration \cite{marspreregistration}\\
$\rightarrow$ Vivianite as ribosomal
  phosphate shadow\\
$\rightarrow$ Planetary cooling argument\\
$\rightarrow$ \textbf{Mars first.
  Panspermia prior inverted.}
\end{quote}

This derivation contains no prior knowledge
of astrobiology at any step.
It arrived here because the geometry of
biological identity, followed to its
foundations, has only one destination:

The LECA attractor is universal.
Carbon chemistry in the right conditions
always produces it.
Mars had those conditions first.
The oldest known biosphere in the solar
system may not be on Earth.

% ── SECTION 9 ─────────────────────────────────────────────────
\section{Conclusion}

Mars was habitable before Earth.
The Cheyava Falls mineral assemblage is
consistent with a mature LUCA-grade
biosphere in terminal collapse at a time
coincident with or prior to the earliest
Earth biosignatures.
The panspermia hypothesis is no longer
exotic --- it is the default hypothesis
requiring elimination given confirmed
Martian biological priority and confirmed
Mars-to-Earth meteorite transfer.

Three predictions are pre-registered
by this document, locked 2026-03-13,
prior to Mars sample return and prior
to any genomic analysis of Martian
material:

\begin{enumerate}[leftmargin=*, label=(\arabic*)]
\item Martian biological activity dates to
  $\geq$3.8\,Ga, at or before the earliest
  confirmed Earth biosignatures.
\item Genomic analysis of Mars sample return
  material will reveal either common origin
  with Earth life or independent biochemical
  architecture.
  There is no third option.
\item The vivianite horizon at Cheyava Falls
  marks the biological cessation boundary
  of the Martian surface biosphere.
  The layer immediately above will be
  abiotic.
\end{enumerate}

The geometry required these conclusions
before the samples come home.

This document is the lock.

\vspace{1em}
\noindent\textcolor{rulecolor}{\rule{\linewidth}{0.4pt}}

% ── ACKNOWLEDGMENTS ───────────────────────────────────────────
\section*{Acknowledgments}

This work was conducted independently without
institutional funding or affiliation.
The derivation sequence is documented in the
pre-registration chain
(DOIs: 10.5281/zenodo.18986790,
10.5281/zenodo.18989573,
10.5281/zenodo.18988844,
10.5281/zenodo.18994082,
10.5281/zenodo.18995862)
and the attractor-oncology repository
(\href{https://github.com/Eric-Robert-Lawson/attractor-oncology}
{github.com/Eric-Robert-Lawson/attractor-oncology}).
The Cheyava Falls data is the work of the NASA
Perseverance mission science team
(Hurowitz et al., \textit{Nature}, 2025)
and is referenced as independent
confirmatory data, not as a derivation source.

% ── REFERENCES ─────────────────────────────────────���──────────
\begin{thebibliography}{99}

\bibitem{preregistration}
Lawson, E.R.\ (2026).
\textit{LECA Resurrection Protocol ---
Pre-Registration v1.0}.
Zenodo.
\href{https://doi.org/10.5281/zenodo.18986790}
{doi:10.5281/zenodo.18986790}

\bibitem{amendmentA1}
Lawson, E.R.\ (2026).
\textit{LECA Resurrection Protocol ---
Registered Amendment A1}.
Zenodo.
\href{https://doi.org/10.5281/zenodo.18989573}
{doi:10.5281/zenodo.18989573}

\bibitem{lucapaper}
Lawson, E.R.\ (2026).
\textit{The Ribosome Is the LUCA}.
Zenodo.
\href{https://doi.org/10.5281/zenodo.18988844}
{doi:10.5281/zenodo.18988844}

\bibitem{uldstandard}
Lawson, E.R.\ (2026).
\textit{The Universal Life Detection Standard}.
Zenodo.
\href{https://doi.org/10.5281/zenodo.18994082}
{doi:10.5281/zenodo.18994082}

\bibitem{marspreregistration}
Lawson, E.R.\ (2026).
\textit{Mars Biosphere Collapse vs.\ Emergence}.
Zenodo.
\href{https://doi.org/10.5281/zenodo.18995862}
{doi:10.5281/zenodo.18995862}

\bibitem{hurowitz2025}
Hurowitz, J.A., Tice, M.M., Allwood, A.C.,
et al.\ (2025).
Redox-driven mineral and organic associations
in Jezero Crater, Mars.
\textit{Nature}, 645(8080), 332--340.
\href{https://doi.org/10.1038/s41586-025-09413-0}
{doi:10.1038/s41586-025-09413-0}

\bibitem{adcock2013}
Adcock, C.T., Hausrath, E.M.,
\& Forster, P.M.\ (2013).
Readily available phosphate from minerals
in early aqueous environments on Mars.
\textit{Nature Geoscience}, 6, 824--827.
\href{https://doi.org/10.1038/ngeo1923}
{doi:10.1038/ngeo1923}

\bibitem{weiss2016}
Weiss, M.C., Sousa, F.L., Mrnjavac, N.,
Neukirchen, S., Roettger, M.,
Nelson-Sathi, S., \& Martin, W.F.\ (2016).
The physiology and habitat of the last
universal common ancestor.
\textit{Nature Microbiology}, 1, 16116.
\href{https://doi.org/10.1038/nmicrobiol.2016.116}
{doi:10.1038/nmicrobiol.2016.116}

\bibitem{wilde2001}
Wilde, S.A., Valley, J.W., Peck, W.H.,
\& Graham, C.M.\ (2001).
Evidence from detrital zircons for the
existence of continental crust and oceans
on the Earth 4.4\,Gyr ago.
\textit{Nature}, 409, 175--178.
\href{https://doi.org/10.1038/35051550}
{doi:10.1038/35051550}

\bibitem{margulis1967}
Sagan, L.\ (1967).
On the origin of mitosing cells.
\textit{Journal of Theoretical Biology},
14(3), 225--274.
\href{https://doi.org/10.1016/0022-5193(67)90079-3}
{doi:10.1016/0022-5193(67)90079-3}

\bibitem{lane2015}
Lane, N.\ (2015).
\textit{The Vital Question}.
W.W.\ Norton.
ISBN: 978-0-393-08881-6

\end{thebibliography}

% ── FOOTER ────────────────────────────────────────────────────
\vspace{1em}
\noindent\textcolor{rulecolor}{\rule{\linewidth}{0.4pt}}
\begin{center}
\footnotesize
\textit{The geometry came first.\ \
Mars was first.\ \
The samples will say the rest.}\\[0.15em]
\href{https://github.com/Eric-Robert-Lawson/attractor-oncology}
{\texttt{github.com/Eric-Robert-Lawson/attractor-oncology}}
\ $\cdot$\
\href{https://doi.org/10.5281/zenodo.18986790}
{\texttt{doi.org/10.5281/zenodo.18986790}}
\ $\cdot$\
\href{https://orcid.org/0009-0002-0414-6544}
{\texttt{ORCID: 0009-0002-0414-6544}}
\end{center}

\end{document}
